\documentclass[11pt]{article}
\usepackage[T1]{fontenc}
\usepackage{lmodern}
\usepackage[margin=1in]{geometry}
\usepackage{amsmath,amssymb,amsthm,mathtools}
\usepackage{microtype,booktabs,longtable,array}
\usepackage{xurl}
\usepackage[hidelinks]{hyperref}
\usepackage{fancyhdr}
\setlength{\headheight}{14pt}
\pagestyle{fancy}\fancyhf{}
\fancyhead[L]{\small Higher spectral defects}\fancyhead[R]{\small ProveIt research continuation}
\fancyfoot[C]{\thepage}
\setlength{\parskip}{4pt}
\setlength{\parindent}{0pt}
\setlength{\emergencystretch}{2em}
\allowdisplaybreaks[2]
\numberwithin{equation}{section}
\newtheorem{theorem}{Theorem}[section]
\newtheorem{proposition}[theorem]{Proposition}
\newtheorem{lemma}[theorem]{Lemma}
\newtheorem{corollary}[theorem]{Corollary}
\theoremstyle{definition}\newtheorem{definition}[theorem]{Definition}
\theoremstyle{remark}\newtheorem{remark}[theorem]{Remark}
\DeclareMathOperator{\Li}{Li}
\DeclareMathOperator{\Rea}{Re}
\DeclareMathOperator{\Ima}{Im}
\DeclareMathOperator{\Log}{Log}
\newcommand{\dd}{\,\mathrm d}
\newcommand{\C}{\mathbb C}
\newcommand{\N}{\mathbb N}
\newcommand{\CC}{\mathcal C}
\newcommand{\KK}{\mathcal K}
\newcommand{\PP}{\mathcal P}
\newcommand{\src}[1]{\nolinkurl{#1}}
\newcommand{\pin}{\src{5a5c0ed6573d25fcabf93653d72449f76df2df79}}
\title{\textbf{Beta Polarization of Higher Spectral Defects}\\[7pt]
\large All-order Stieltjes counterterms, Gamma integrals,\\
polylogarithmic second jets, and confluent quadratic masses}
\author{Research continuation prepared for Vladimir Reshetnikov\\
\small ProveIt / Analysis / Polylogarithms\\
\small Mathematical derivations and computational checks with OpenAI assistance}
\date{11 October 2026}
\hypersetup{pdftitle={Beta Polarization of Higher Spectral Defects},pdfauthor={Research continuation prepared for Vladimir Reshetnikov},pdfsubject={Exact Hurwitz, Gamma, harmonic, Stieltjes and polylogarithmic identities}}
\begin{document}
\maketitle
\begin{abstract}
We resolve the explicit higher-defect and two-mass second-derivative
questions posed in the quadratic-spectrum continuation of the ProveIt
programme. An endpoint-subtracted Euler beta integral gives every
spectral defect of order $r$ as a finite combination of ordinary
logarithmically weighted integrals of linear Hurwitz-zeta derivatives of
orders at most $r-1$. A separate, normally convergent coefficient formula
makes all resonant generalized Stieltjes terms explicit, with coefficients
in generalized harmonic numbers.

At second order the nonlocal part is a Gamma--Hurwitz chord integral and
the correction is a finite odd-weight polynomial. In contrast to the
first-derivative parity theorem, no nonzero polynomial weight annihilates
the second defect for every mass. The scalar case gives convergent
arctangent-square, log-Gamma, harmonic-zeta, and polygamma identities.
For every nonpositive integral spectral argument, the bilateral second
jet is expressed by finitely many complex-shift Hurwitz second derivatives
and one ordinary integral of an explicit finite polylogarithm expression.

For any number of quadratic factors with an even polynomial weight, the
second jet has a pairwise covariance formula. Each covariance is a
squared-mass beta chord of the already evaluated first jet; thus it has
finite Gamma or polylogarithmic integrands. Its full confluent expansion
has coefficients $2(-1)^kH_{k-1}Z_W(k)/k$, with finite polygamma reductions.
All parameter primitives are anchored explicitly. The article includes
proofs, executable exact checks, numerical diagnostics, an integration
fragment, and a scoped provenance audit. The beta identity and
first-derivative anomaly mechanism are classical inputs, not priority
claims; minimal-function-algebra questions remain separate.
\end{abstract}
\medskip
\noindent\textbf{Status.} Ordinary analytic proofs with reproducible
finite checks; not proof-assistant formalized or independently peer reviewed.
No repository files were modified.
\clearpage
\tableofcontents
\clearpage

\section{The research question and the contribution}\label{hsd:sec:scope}
The canonical ProveIt manuscript distinguishes exact analytic identities,
relations in a specified formal presentation, and numerical residuals
\cite{hsd:repo}. This distinction is particularly important for
zeta-regularized logarithms: a pointwise logarithmic identity need not
survive spectral continuation without a correction.

The immediate predecessor, \emph{Quadratic-Spectrum Stieltjes Identities}
\cite{hsd:predecessor}, proves a first-derivative factorization formula and
asks in its Section~7.6 for three further developments. Question~2 asks
for an explicit separation of the nonlocal and residue parts of higher
factorization defects. Question~7 asks for the second derivative of a
product of two quadratic spectra, including its coincident-mass limit,
without an unspecified subtraction of divergent sums. Question~1 asks
for useful function coordinates for higher derivatives at negative
integers, and separately whether a smaller finite function algebra suffices.

Theorems~\ref{hsd:thm:beta-jets} and \ref{hsd:thm:stieltjes}
give all-order formulas for the first problem.
Theorems~\ref{hsd:thm:gamma-second} and \ref{hsd:thm:nonvanishing}
give a particularly explicit second-order answer and its exact universal
vanishing classification. Theorems~\ref{hsd:thm:masses} and
\ref{hsd:thm:collision} solve the two-mass representation and confluence
problem, and also handle arbitrarily many masses.
Theorem~\ref{hsd:thm:negative-second} supplies a finite-integrand
representation at every negative integer. It does \emph{not} prove that
the remaining ordinary integral is irreducible to a smaller algebra.

The main mechanism is order lowering. The subtracted beta kernel has a
simple zero in the spectral parameter at zero. Consequently an $r$th
defect needs linear Hurwitz derivatives only through order $r-1$ inside
the integral. At order two these are first Hurwitz derivatives, including
$\log\Gamma$. An independent inverse-power expansion identifies exactly
where generalized Stieltjes constants replace ordinary zeta values.

\subsection{Classical inputs and the priority boundary}
Euler's beta integral, its two-denominator parameterization, Gamma
products, Hurwitz continuation and differentiation, and elementary
polylogarithm relations are classical \cite{hsd:beta,hsd:gamma,hsd:hurwitz,hsd:polylog}.
The multiplicative anomaly for first derivatives of spectral zeta
functions is also an established subject; Dowker explicitly treats
products of linear and second-order factors \cite{hsd:dowker}.
The predecessor's finite first-jet polylogarithm formulas are retained
with attribution, not relabeled as new results.

The contribution here is the precise endpoint-subtracted all-jet
calculus, its harmonic/Stieltjes coefficient implementation, the second
universal-vanishing theorem, and the explicit confluent and primitive
consequences. No comprehensive historical-priority or arithmetic
independence claim is made. Neither the Gaussian $S_6$ nor the revised
$S_8$ special-value problem is solved by this article.

\subsection{Source identity and scope of inspection}
The predecessor's single-file TeX was read in full for the relevant
sections and materialized for comparison. Its SHA-256 is
\begin{center}\small
\src{4185cc4aa2299c8f417ca376b3ddc4af3fab58752d4cd93f70c2afad661f95aa}.
\end{center}
This agrees exactly with the repository's sixteenth-intake manifest at
\pin. The recorded path is
\begin{quote}\small
\src{Analysis/Polylogarithms/docs/reports/stieltjes-correlation-calculus/continuations/quadratic-spectrum-stieltjes/article_singlefile.tex}.
\end{quote}
Thus the mathematical target is a byte-identified repository report,
not a conjecture inferred from a filename. The canonical preface,
differentiation and research-status sections, incoming catalogue, and
relevant intake metadata were also inspected. This was a targeted audit,
not an exhaustive reading of all incoming reports or the entire repository.
Detailed source records are included in \src{notes/SOURCE_AUDIT.md}.

\section{Analytic normalization and an endpoint-subtracted beta master}
\label{hsd:sec:beta}
Fix $a>0$, write $x_n=n+a$, and let
\begin{equation}\label{hsd:eq:W}
 W(x)=\sum_{d=0}^D w_dx^d\in\C[x].
\end{equation}
Coefficients outside this range are zero. All spectral superscripts
$(r)$ below differentiate the spectral variable, not the shift.
On the right half-plane, $\Log$ and $\log\Gamma$ mean their holomorphic
branches that are real on the positive axis. In particular,
$\log\Gamma(z)$ is not formed by taking the principal logarithm of a
numerically evaluated $\Gamma(z)$.

For $\Rea(a+z)>0$, define initially for $\Rea s>D+1$
\begin{equation}\label{hsd:eq:F}
 F_{a,W}(s;z)=\sum_{n\ge0}W(x_n)(x_n+z)^{-s}.
\end{equation}
A finite binomial expansion gives its continuation:
\begin{equation}\label{hsd:eq:Ffinite}
 F_{a,W}(s;z)=\sum_{d=0}^D w_d\sum_{h=0}^d
 \binom dh(-z)^{d-h}\zeta(s-h,a+z).
\end{equation}
It is holomorphic at $s=0$, locally uniformly in $z$ in the stated domain.
Put
\begin{equation}\label{hsd:eq:phi}
 \Phi_{a,W}(z)=F_{a,W}'(0;z),\qquad
 A_W(z)=\sum_{d=0}^D\frac{(-1)^{d+1}w_dz^{d+1}}{d+1}.
\end{equation}
For $W=1$ the classical Lerch formula gives
$\Phi_{a,1}(z)=\log\Gamma(a+z)-\tfrac12\log(2\pi)$.

For endpoints $u,v$ with $\Rea(a+u),\Rea(a+v)>0$, define
\begin{equation}\label{hsd:eq:Zuv}
 Z_{a,W}^{u,v}(s)=\sum_{n\ge0}W(x_n)
 \exp\{-s(\Log(x_n+u)+\Log(x_n+v))\}.
\end{equation}
The definition uses the \emph{sum} of the two specified logarithms.
Its initial convergence domain is $2\Rea s>D+1$.
Set
\begin{align}
 G_{a,W}(s;u,v)&=\tfrac12\{F_{a,W}(2s;u)+F_{a,W}(2s;v)\}
                    -Z_{a,W}^{u,v}(s),\label{hsd:eq:G}\\
 \Delta_r(a,W;u,v)&=G_{a,W}^{(r)}(0;u,v).\label{hsd:eq:Delta}
\end{align}
For $u=ib,v=-ib$, write $Z_W(s;a,b)$ and $\Delta_r(W;a,b)$.
Then $Z_W(s;a,b)=\sum W(x_n)(x_n^2+b^2)^{-s}$, and
\begin{equation}\label{hsd:eq:Delta-conjugate}
 \Delta_r(W;a,b)=2^{r-1}\bigl(F_{a,W}^{(r)}(0;ib)
 +F_{a,W}^{(r)}(0;-ib)\bigr)-Z_W^{(r)}(0;a,b).
\end{equation}

For an analytic function $f$ on a neighborhood of the segment $[u,v]$,
define the symmetric endpoint subtraction
\begin{equation}\label{hsd:eq:C}
 \CC_t f(u,v)=f(u)+f(v)-f((1-t)u+tv)-f(tu+(1-t)v).
\end{equation}
This vanishes to first order at each endpoint of $0\le t\le1$.
More precisely,
\begin{equation}\label{hsd:eq:C-limit}
 \lim_{t\to0,1}\frac{\CC_t f(u,v)}{t(1-t)}
 =(v-u)\bigl(f'(v)-f'(u)\bigr).
\end{equation}
The limit also holds at $u=v$, where both sides are zero.

\begin{lemma}[The zero jet is a local polynomial]\label{hsd:lem:zero}
For all admissible $a,z,W$,
\begin{equation}\label{hsd:eq:Fzero}
 F_{a,W}(0;z)=\sum_{d=0}^D w_d\zeta(-d,a)+A_W(z).
\end{equation}
In particular, $\CC_tF_{a,W}(0;\cdot)=\CC_t A_W$ is independent of $a$.
\end{lemma}
\begin{proof}
For $|z|<a$, expand the summand binomially and continue:
\begin{equation}\label{hsd:eq:Fbinomial}
 F_{a,W}(s;z)=\sum_{d=0}^D w_d\sum_{k\ge0}
 \frac{(-1)^k(s)_kz^k}{k!}\zeta(s+k-d,a).
\end{equation}
The $k=0$ term gives $\zeta(-d,a)$. Every $k\ge1$ term has a factor
$s$ and vanishes at zero except $k=d+1$, where the simple Hurwitz pole
cancels it. Since $(s)_k/k!=s/k+O(s^2)$, that surviving term is
$(-1)^kz^k/k$. The series and its pole-subtracted terms converge locally
normally for $|z|<a$; the ratio $|z|/a<1$ dominates the polynomial
coefficient growth. This proves the formula locally. Both sides are
holomorphic on $\Rea(a+z)>0$, so analytic continuation proves it there.
\end{proof}

\begin{theorem}[Subtracted beta master]\label{hsd:thm:beta}
The function $Z_{a,W}^{u,v}$ has a continuation holomorphic at zero.
As an identity of germs at $s=0$,
\begin{equation}\label{hsd:eq:beta-master}
 G_{a,W}(s;u,v)=\frac{s}{4}B_*(s)
 \int_0^1\frac{[t(1-t)]^s}{t(1-t)}
       \CC_t F_{a,W}(2s;\cdot)(u,v)\dd t,
\end{equation}
where
\begin{equation}\label{hsd:eq:Bstar}
 B_*(s)=\frac{\Gamma(1+2s)}{\Gamma(1+s)^2}.
\end{equation}
The integral in \eqref{hsd:eq:beta-master} is an ordinary convergent
integral near zero, jointly holomorphic in $s,u,v$ in their stated local
domains. Every fixed spectral derivative may be passed under it.
\end{theorem}
\begin{proof}
For $\Rea s>(D+1)/2$ use the classical two-denominator beta identity,
which follows from Euler's beta integral or two Gamma integrals:
\[
 A^{-s}B^{-s}=\frac{\Gamma(2s)}{\Gamma(s)^2}
 \int_0^1\frac{t^{s-1}(1-t)^{s-1}}{((1-t)A+tB)^{2s}}\dd t.
\]
It holds for $\Rea A,\Rea B>0$ with the prescribed logarithms.
The absolute convergence of the $n$-sum is uniform in $t$, apart from
the integrable beta weight. Sum it with
$A=x_n+u$, $B=x_n+v$, and average the result under $t\mapsto1-t$.
Since the beta weight integrates to $\Gamma(s)^2/\Gamma(2s)$, subtracting
the endpoint average produces
\[
 G(s)=\frac{\Gamma(2s)}{2\Gamma(s)^2}
       \int_0^1t^{s-1}(1-t)^{s-1}\CC_tF(2s;\cdot)\dd t.
\]
The prefactor equals $sB_*(s)/4$.

For a small closed disk about $s=0$, the finite formula
\eqref{hsd:eq:Ffinite} is holomorphic, uniformly on the compact endpoint
segments under consideration. Taylor expansion in $t$ gives
$\CC_tF(2s;\cdot)=O(t)$ at zero and $O(1-t)$ at one, uniformly on that
disk. Each fixed $s$-derivative adds only a fixed power of
$|\log(t(1-t))|$. Thus the integral and all such derivatives are
absolutely convergent and locally dominated whenever $\Rea s>-1$
in this small disk. This proves holomorphy and differentiation under the
integral.

The right side continues meromorphically from the initial half-plane,
with the isolated positive spectral poles of the finite Hurwitz expression
bypassed. It is holomorphic at zero. Defining $Z$ by
\eqref{hsd:eq:G} there proves both its continuation and the identity.
The argument is locally uniform in $u,v$ and covers coincident endpoints.
\end{proof}

\subsection{All defects as finite lower-jet integrals}
Write
\begin{equation}\label{hsd:eq:beta-coeffs}
 B_*(s)=\sum_{m\ge0}\beta_ms^m
 =\exp\left(\sum_{n\ge2}
     \frac{(-1)^n(2^n-2)\zeta(n)}{n}s^n\right).
\end{equation}
Hence $\beta_0=1$, $\beta_1=0$, $\beta_2=\zeta(2)$,
$\beta_3=-2\zeta(3)$, and
$\beta_4=\tfrac72\zeta(4)+\tfrac12\zeta(2)^2$.
An exact finite recurrence is
\begin{equation}\label{hsd:eq:beta-rec}
 m\beta_m=\sum_{n=2}^m(-1)^n(2^n-2)\zeta(n)\beta_{m-n}.
\end{equation}
Put $\lambda(t)=\log(t(1-t))$ and
\begin{equation}\label{hsd:eq:I}
 I_{j,p}(a,W;u,v)=\int_0^1\frac{\lambda(t)^p}{t(1-t)}
 \CC_t\bigl(F_{a,W}^{(j)}(0;\cdot)\bigr)(u,v)\dd t.
\end{equation}
No finite-part convention is involved in this integral.

\begin{theorem}[All-order lower-jet reduction]\label{hsd:thm:beta-jets}
For every integer $r\ge1$ and every admissible endpoint pair,
\begin{equation}\label{hsd:eq:all-beta}
 \boxed{\displaystyle
 \Delta_r(a,W;u,v)=\frac{r!}{4}
 \sum_{m+j+p=r-1}\frac{2^j\beta_m}{j!p!}I_{j,p}(a,W;u,v).}
\end{equation}
All sums are finite. The maximal spectral derivative inside an integral
is $r-1$. Every $j=0$ term is an explicit polynomial in $u,v$ independent
of $a$.
\end{theorem}
\begin{proof}
Expand $B_*(s)$, $\exp(s\lambda(t))$ and
$F(2s;z)=\sum_{j\ge0}2^js^jF^{(j)}(0;z)/j!$ in
\eqref{hsd:eq:beta-master}. The domination established in
Theorem~\ref{hsd:thm:beta} justifies extraction of each coefficient under
the integral. The outer factor $s/4$ gives exactly the constraint
$m+j+p=r-1$. Lemma~\ref{hsd:lem:zero} identifies the local terms.
\end{proof}

The first four rows, suppressing common arguments, are
\begin{align}
 \Delta_1&=\tfrac14 I_{0,0},\label{hsd:eq:beta-row1}\\
 \Delta_2&=I_{1,0}+\tfrac12I_{0,1},\label{hsd:eq:beta-row2}\\
 \Delta_3&=3I_{2,0}+3I_{1,1}+\tfrac34I_{0,2}
                         +\tfrac32\zeta(2)I_{0,0},\label{hsd:eq:beta-row3}\\
 \Delta_4&=8I_{3,0}+12I_{2,1}+6I_{1,2}+I_{0,3}
       +12\zeta(2)I_{1,0}+6\zeta(2)I_{0,1}
       -12\zeta(3)I_{0,0}.\label{hsd:eq:beta-row4}
\end{align}
At conjugate endpoints and for even $W$, the polynomial $A_W$ is odd,
so $\CC_tA_W(ib,-ib)=0$ identically. All $I_{0,p}$ then vanish.
This removes the local polynomial terms but does not remove the
positive-$j$ integrals. That distinction is the reason the
first-derivative parity theorem does not extend unchanged to higher jets.

\section{A global harmonic--Stieltjes coefficient formula}
\label{hsd:sec:stieltjes}
The beta master is finite in the number of function coordinates at every
order. A different representation is useful for exact coefficients and
for evaluating resonances without subtracting two large meromorphic terms.
Use the convention
\begin{equation}\label{hsd:eq:Stieltjes}
 \zeta(1+\varepsilon,X)=\frac1\varepsilon+
 \sum_{n\ge0}\frac{(-1)^n\gamma_n(X)}{n!}\varepsilon^n,
 \qquad \gamma_0(X)=-\psi(X).
\end{equation}
Let
\begin{equation}\label{hsd:eq:oddH}
 O_k^{(p)}=\sum_{j=0}^{k-1}\frac1{(2j+1)^p}
 =H_{2k}^{(p)}-2^{-p}H_k^{(p)},\qquad O_k=O_k^{(1)}.
\end{equation}
The coefficient polynomial controlling the conjugate defect is
\begin{align}
 D_k(s)&=\frac{(2s)_{2k}}{(2k)!}-\frac{(s)_k}{k!}
       =\sum_{\ell\ge2}d_{k,\ell}s^\ell\label{hsd:eq:D}\\
 &=\frac{s}{k}\prod_{j=1}^{k-1}\left(1+\frac{s}{j}\right)
 \left\{\prod_{j=0}^{k-1}\left(1+\frac{2s}{2j+1}\right)-1\right\}.
 \label{hsd:eq:Dprod}
\end{align}
The second equality follows from duplication of the rising factorial.
It exhibits the double zero without cancellation in floating-point
arithmetic. In particular,
\begin{align}
 d_{k,2}&=\frac{2O_k}{k},\label{hsd:eq:d2}\\
 d_{k,3}&=\frac2k\bigl(O_kH_{k-1}+O_k^2-O_k^{(2)}\bigr).
 \label{hsd:eq:d3}
\end{align}
Every higher coefficient is a finite polynomial in generalized harmonic
numbers. One can either multiply the finite products in
\eqref{hsd:eq:Dprod}, or use
\[
 \log\prod_{j=1}^{k-1}(1+s/j)
 =\sum_{p\ge1}\frac{(-1)^{p-1}}pH_{k-1}^{(p)}s^p,
 \qquad
 \log\prod_{j=0}^{k-1}(1+2s/(2j+1))
 =\sum_{p\ge1}\frac{(-1)^{p-1}2^p}pO_k^{(p)}s^p.
\]
Only finitely many coefficients are needed for a prescribed jet.

For $r\ge1$ define the elementary finite-head term
\begin{equation}\label{hsd:eq:qr}
 q_r(x,b)=2^{r-1}\{(-\Log(x+ib))^r+(-\Log(x-ib))^r\}
                   -(-\log(x^2+b^2))^r.
\end{equation}
Thus $q_1=0$ and
\begin{equation}\label{hsd:eq:q2}
 q_2(x,b)=\bigl(\Log(x+ib)-\Log(x-ib)\bigr)^2
        =-4\arctan^2(b/x).
\end{equation}

\begin{theorem}[Every resonant coefficient, with a global finite head]
\label{hsd:thm:stieltjes}
Let $a,b>0$, $r\ge1$, and choose any integer $N\ge0$ such that
$X=a+N>b$. For an integer $\tau$ put
\begin{equation}\label{hsd:eq:Tnonres}
 T_{r,k}(\tau;X)=r!\sum_{\ell=2}^{r}
 \frac{2^{r-\ell}d_{k,\ell}}{(r-\ell)!}
             \zeta^{(r-\ell)}(\tau,X),\qquad \tau\ne1,
\end{equation}
and at the resonance set
\begin{equation}\label{hsd:eq:Tres}
 T_{r,k}(1;X)=r!\left\{\frac{d_{k,r+1}}2+
 \sum_{\ell=2}^{r}\frac{(-2)^{r-\ell}d_{k,\ell}}{(r-\ell)!}
                         \gamma_{r-\ell}(X)\right\}.
\end{equation}
Empty sums are zero. Then
\begin{equation}\label{hsd:eq:Stieltjes-master}
 \boxed{\displaystyle
 \Delta_r(W;a,b)=\sum_{n=0}^{N-1}W(a+n)q_r(a+n,b)
 +\sum_{d=0}^D w_d\sum_{k\ge1}(-1)^kb^{2k}
                         T_{r,k}(2k-d;X).}
\end{equation}
The double sum is absolutely and locally normally convergent. It is
independent of the permissible cutoff $N$. The largest Stieltjes index
which can occur is $r-2$.
\end{theorem}
\begin{proof}
On the tail $x\ge X>b$, expand $\tfrac12\{(x+ib)^{-2s}+
(x-ib)^{-2s}\}$ and $(x^2+b^2)^{-s}$ in powers of $b/x$.
Odd powers cancel in the first expression. Subtracting gives
\begin{align}\label{hsd:eq:Gtail}
 G(s;ib,-ib)={}&\sum_{n<N}\tfrac12 W(x_n)
 \{(x_n+ib)^{-2s}+(x_n-ib)^{-2s}-2(x_n^2+b^2)^{-s}\}\notag\\
 &+\sum_{d=0}^D\sum_{k\ge1} w_d(-1)^kb^{2k}
                         D_k(s)\zeta(2s+2k-d,X),
\end{align}
where $0\le d\le D$ and $k\ge1$.
For $k$ sufficiently large, the Hurwitz argument has large positive
real part. Uniformly for $s$ on a fixed small circle,
$\zeta(2s+2k-d,X)$ is bounded by a constant times
$X^{-2k+d-2\Rea s}$, up to a harmless polynomial factor. The rising
factorial ratios grow only polynomially in $k$ on such a circle.
Therefore $(b/X)^{2k}$ dominates the tail normally. A fixed number of
spectral or shift derivatives introduces at most polynomial and
logarithmic factors. The finitely many possible Hurwitz poles are
handled individually.

If $2k-d\ne1$, Taylor multiplication gives
\eqref{hsd:eq:Tnonres}. At $2k-d=1$, insert
\eqref{hsd:eq:Stieltjes} with $\varepsilon=2s$. Because $D_k$ starts
at $s^2$, the product is holomorphic; its $r$th derivative is exactly
\eqref{hsd:eq:Tres}. The finite head differentiates to $Wq_r$.
This proves the formula and all convergence assertions. Different
cutoffs give the same continued function before differentiation, proving
cutoff independence.
\end{proof}

For clarity, the first three resonant rows are
\begin{align}
 T_{1,k}(1;X)&=\tfrac12d_{k,2},\label{hsd:eq:res1}\\
 T_{2,k}(1;X)&=d_{k,3}+2d_{k,2}\gamma_0(X),\label{hsd:eq:res2}\\
 T_{3,k}(1;X)&=3d_{k,4}+6d_{k,3}\gamma_0(X)
                         -12d_{k,2}\gamma_1(X).\label{hsd:eq:res3}
\end{align}
The term $d_{k,r+1}/2$ is easily lost by truncating $D_k$ at order $r$
before multiplying by the pole. Such a truncation gives an incorrect
jet even though $D_k(0)=D_k'(0)=0$.

At $r=1$, only $d=2k-1$ contributes, and
\eqref{hsd:eq:Stieltjes-master} recovers the predecessor's identity
\begin{equation}\label{hsd:eq:first-anomaly}
 \Delta_1(W;a,b)=\sum_{k\ge1}
             \frac{(-1)^kw_{2k-1}b^{2k}}k O_k.
\end{equation}
This is inherited, not a new first-order theorem. At higher orders,
the same finite resonance set remains, but infinitely many nonresonant
coordinates survive. The resonant Stieltjes values in
\eqref{hsd:eq:Tres} should not be confused with the shift-independent
polynomial part of the beta-integral formula: the two formulas organize
the same continued expression in different bases.

\section{The exact second defect and its vanishing classification}
\label{hsd:sec:gamma}
Define the Gamma--Hurwitz polarization
\begin{equation}\label{hsd:eq:K}
 \KK_{a,W}(u,v)=\int_0^1\frac{\CC_t\Phi_{a,W}(u,v)}{t(1-t)}\dd t.
\end{equation}
For a fixed polynomial weight it is one ordinary integral of the finite
Hurwitz expression \eqref{hsd:eq:Ffinite}. For $W=1$ it involves only
$\log\Gamma$. Its integrand has removable endpoints by
\eqref{hsd:eq:C-limit}; it is jointly holomorphic and vanishes when $u=v$.

\begin{theorem}[Gamma integral plus an explicit odd-weight polynomial]
\label{hsd:thm:gamma-second}
For every polynomial $W$ and every $a,b>0$,
\begin{equation}\label{hsd:eq:second-main}
 \boxed{\displaystyle
 \Delta_2(W;a,b)=\KK_{a,W}(ib,-ib)-C_W(b),\qquad
 C_W(b)=\sum_{k\ge1}\frac{2(-1)^kw_{2k-1}b^{2k}}k
                         \bigl(O_k^2+O_k^{(2)}\bigr).}
\end{equation}
The correction is finite and independent of $a$. In particular,
\begin{equation}\label{hsd:eq:quadratic-second-finite}
 Z_W''(0;a,b)=2\{F_{a,W}''(0;ib)+F_{a,W}''(0;-ib)\}
                      -\KK_{a,W}(ib,-ib)+C_W(b).
\end{equation}
There is no undefined regularized sum on the right side.
\end{theorem}
\begin{proof}
The second beta row gives
$\Delta_2=\KK+\tfrac12 I_{0,1}$. For conjugate endpoints,
Lemma~\ref{hsd:lem:zero} gives
\[
 \CC_tA_W(ib,-ib)=\sum_{k\ge1}
       \frac{(-1)^kw_{2k-1}b^{2k}}k\bigl(1-(1-2t)^{2k}\bigr).
\]
We evaluate the logarithmic moment, including its normalization.
Put
\[
 P_k(s)=\prod_{j=0}^{k-1}\left(1+\frac{2s}{2j+1}\right),\qquad
 J_k(s)=\int_0^1[t(1-t)]^{s-1}\bigl(1-(1-2t)^{2k}\bigr)\dd t.
\]
By $y=1-2t$ and Euler's beta integral,
\begin{equation}\label{hsd:eq:Jk}
 J_k(s)=\frac{2}{sB_*(s)}\bigl(1-P_k(s)^{-1}\bigr).
\end{equation}
This is first obtained for $\Rea s>0$ and is holomorphic at zero;
its defining subtracted integral is locally integrable for $\Rea s>-1$.
Since
\[
 P_k(s)^{-1}=1-2O_ks+2\bigl(O_k^2+O_k^{(2)}\bigr)s^2+O(s^3),
 \qquad B_*(s)=1+O(s^2),
\]
we have $J_k(0)=4O_k$ and
$J_k'(0)=-4(O_k^2+O_k^{(2)})$. Thus
$\tfrac12I_{0,1}=-C_W(b)$. The second formula follows directly from the
definition of $\Delta_2$.
\end{proof}

The first correction terms in $\Delta_2-\KK$ are
\begin{equation}\label{hsd:eq:local-examples}
 4w_1b^2-\frac{26}{9}w_3b^4+\frac{1576}{675}w_5b^6
                         -\frac{21946}{11025}w_7b^8+\cdots.
\end{equation}
For real weights, an equivalent real integral is
\begin{equation}\label{hsd:eq:Kreal}
 \KK_{a,W}(ib,-ib)=8\int_0^1
 \frac{\Rea\Phi_{a,W}(ib)-\Rea\Phi_{a,W}(ibt)}{1-t^2}\dd t.
\end{equation}
For complex weights, use \eqref{hsd:eq:K} rather than pulling coefficients
through a real-part operation.

\subsection{An independent coefficient derivation}
A useful consistency check on the beta proof is the Taylor expansion
\begin{equation}\label{hsd:eq:PhiTaylor}
 \Phi_{a,W}(z)=\Phi_{a,W}(0)+\sum_{j\ge1}\frac{(-1)^jz^j}{j}
 \left\{\sum_{d=0}^D w_d\widehat\zeta(j-d,a)+w_{j-1}H_{j-1}\right\},
 \quad |z|<a,
\end{equation}
where $\widehat\zeta(1,a)=\gamma_0(a)$ and
$\widehat\zeta(s,a)=\zeta(s,a)$ at the other displayed integer arguments.
To prove it, differentiate \eqref{hsd:eq:Fbinomial}; the resonant term
uses $(s)_j/j!=s(1+H_{j-1}s+O(s^2))/j$.
The harmonic addend is essential.

For conjugate endpoints, direct polynomial integration gives
\[
 \int_0^1\frac{\CC_t(z^{2k})(ib,-ib)}{t(1-t)}\dd t
                    =8(-1)^kO_kb^{2k},
 \qquad \CC_t(z^{2k+1})(ib,-ib)=0.
\]
Inserting \eqref{hsd:eq:PhiTaylor} into $\KK$ and using
$H_{2k-1}=O_k+\tfrac12H_{k-1}$ reproduces exactly
\eqref{hsd:eq:second-main} from the second Stieltjes row.
This explains why an uncorrected termwise ``sum of logarithm squares''
can differ from $\KK$ for an odd polynomial weight.

\begin{theorem}[Universal vanishing at second order]\label{hsd:thm:nonvanishing}
Fix $a>0$ and $W\in\C[x]$. Then
\begin{equation}\label{hsd:eq:zero-classification}
 \Delta_2(W;a,b)=0\ \text{for every }b>0
 \quad\Longleftrightarrow\quad W=0.
\end{equation}
In fact, for a nonzero $W$, all sufficiently high even Taylor
coefficients in $b$ at zero are nonzero.
\end{theorem}
\begin{proof}
For $0<b<a$ use \eqref{hsd:eq:Stieltjes-master} with $N=0$.
When $2k>D+1$, there is no resonance, and the coefficient of $b^{2k}$ is
\begin{equation}\label{hsd:eq:tailmoment}
 \frac{4(-1)^kO_k}{k}\sum_{d=0}^D w_d\zeta(2k-d,a)
 =\frac{4(-1)^kO_k}{k}\sum_{n\ge0}\frac{W(a+n)}{(a+n)^{2k}}.
\end{equation}
Suppose $W\ne0$, and let $n_0$ be the smallest nonnegative integer with
$W(a+n_0)\ne0$. It exists because a nonzero polynomial has finitely many
roots. Put $x_0=a+n_0$. Multiplying the last sum by $x_0^{2k}$ gives
\[
 W(x_0)+\sum_{n>n_0}W(a+n)\left(\frac{x_0}{a+n}\right)^{2k}
                         \longrightarrow W(x_0)\ne0.
\]
To justify the limit for complex coefficients as well, fix $k_0$ with
$2k_0>D+1$. The absolute values for $k\ge k_0$ are dominated by a
convergent polynomially weighted Dirichlet series. Dominated convergence
then applies. Thus the coefficients in \eqref{hsd:eq:tailmoment} are
nonzero for every sufficiently large $k$, contradicting an identically
zero analytic function. The converse is immediate.
\end{proof}

This theorem concerns an identity valid for all masses, not zeros at one
chosen mass. It also does not assert positivity for arbitrary weights:
continued polynomially weighted sums can have either sign. The exact
contrast is that even $W$ annihilates \emph{every first defect}, whereas
only $W=0$ annihilates \emph{every second defect}.

\section{Ordinary Gamma, polygamma, and harmonic identities}
\label{hsd:sec:ordinary}
The scalar specialization needs no regularization at all. It provides
both concrete identities and a direct way to see the nonlocal term.

\begin{proposition}[A convergent logarithm-square polarization]
\label{hsd:prop:ordinaryK}
For $a>0$ and $\Rea(a+u),\Rea(a+v)>0$,
\begin{equation}\label{hsd:eq:ordinaryK}
 \KK_{a,1}(u,v)=\sum_{n\ge0}
       \bigl(\Log(n+a+u)-\Log(n+a+v)\bigr)^2.
\end{equation}
The series is absolutely and locally normally convergent. The square is
an ordinary complex square, not an absolute square. Its mixed derivative is
\begin{equation}\label{hsd:eq:Kdigamma}
 \partial_u\partial_v\KK_{a,1}(u,v)
 =-2\frac{\psi(a+v)-\psi(a+u)}{v-u},
\end{equation}
with diagonal value $-2\psi'(a+u)$.
\end{proposition}
\begin{proof}
For fixed nonzero $A,B$ in the right half-plane, the scalar beta formula
of Theorem~\ref{hsd:thm:beta}, without any sum, gives on taking second
jets
\[
 \int_0^1\frac{\CC_t(-\Log(x+z))(u,v)}{t(1-t)}\dd t
       =\bigl(\Log(x+u)-\Log(x+v)\bigr)^2.
\]
Equivalently this follows by two elementary differentiations in $u,v$
and evaluation on the diagonal. In the Weierstrass expansion of
$\log\Gamma(a+z)$, the affine terms in $z$ are annihilated by $\CC_t$.
The remaining tail is $O((n+a)^{-2})$, uniformly on compact endpoint
sets, including after the displayed endpoint subtraction. Thus the
scalar identity may be summed under the integral, proving
\eqref{hsd:eq:ordinaryK}. Differentiating the normally convergent series
and using the convergent digamma difference proves
\eqref{hsd:eq:Kdigamma}. The limit at the diagonal follows by continuity.
\end{proof}

The collision expansion of this Gamma identity is particularly simple.
For admissible endpoints and $|\delta|<|a+u|$,
\begin{align}
 \KK_{a,1}(u,u+\delta)
 &=\sum_{k\ge2}\frac{2(-1)^kH_{k-1}}k\delta^k\zeta(k,a+u)
 \label{hsd:eq:scalar-collision-zeta}\\
 &=2\sum_{k\ge2}\frac{H_{k-1}}{k!}\delta^k\psi^{(k-1)}(a+u).
 \label{hsd:eq:scalar-collision-psi}
\end{align}
Indeed,
$\log^2(1+z)=\sum_{k\ge2}2(-1)^kH_{k-1}z^k/k$ for $|z|<1$.
The endpoint hypothesis fixes the logarithm; since $\Rea(a+u)>0$,
$|n+a+u|\ge|a+u|$, which justifies normal summation of this series.
The first terms are
\[
 \delta^2\zeta(2,a+u)-\delta^3\zeta(3,a+u)
                     +\frac{11}{12}\delta^4\zeta(4,a+u)+O(\delta^5).
\]

For $a,b>0$, put $S_a(b)=\sum_{n\ge0}\arctan^2(b/(n+a))$.
\begin{corollary}[Arctangent--Gamma identities]
\label{hsd:cor:atan}
The following identities hold:
\begin{align}
 \Delta_2(1;a,b)&=-4S_a(b),\label{hsd:eq:scalar-defect}\\
 S_a(b)&=2b\int_0^b
 \frac{\Rea\log\Gamma(a+iy)-\Rea\log\Gamma(a+ib)}{b^2-y^2}\dd y,
 \label{hsd:eq:atan-gamma}\\
 S_a(b)&=\sum_{k\ge1}\frac{(-1)^{k-1}O_k}{k}
                         b^{2k}\zeta(2k,a),\qquad 0<b<a.
 \label{hsd:eq:atan-harmonic}
\end{align}
For arbitrary $b$, the global version of the last formula is obtained by
adding $\sum_{n<N}\arctan^2(b/(a+n))$ and replacing $a$ by $a+N>b$
in its zeta tail.
\end{corollary}
\begin{proof}
Apply \eqref{hsd:eq:ordinaryK} to $u=ib,v=-ib$, use
\eqref{hsd:eq:q2}, and then use \eqref{hsd:eq:Kreal} with $y=bt$.
The last identity follows either by the Taylor series of $\arctan^2$
or by \eqref{hsd:eq:Stieltjes-master} for $W=1,r=2$.
\end{proof}

In particular the constant weight, which has zero first defect, has a
strictly negative second defect at every positive mass. This is an exact
counterexample to an unqualified extension of the first-order parity law.
The predecessor already warned against that extension; this corollary
supplies the explicit obstruction rather than correcting a false claim
in that report.

At an integral shift $a=N+1$, the same identities become identities in
ordinary harmonic tails, since
\[
 \zeta(2k,N+1)=\zeta(2k)-H_N^{(2k)}.
\]
For example, on $0<b<N+1$,
\begin{align}
 &2b\int_0^b\frac{\Rea\log\Gamma(N+1+iy)
              -\Rea\log\Gamma(N+1+ib)}{b^2-y^2}\dd y\notag\\
 &\quad=\sum_{k\ge1}\frac{(-1)^{k-1}}k
  \left(H_{2k}-\frac12H_k\right)b^{2k}
                     \bigl(\zeta(2k)-H_N^{(2k)}\bigr).
 \label{hsd:eq:explicit-tail-Gamma}
\end{align}
This is an equality of ordinary convergent quantities with every
coefficient explicit.

\subsection{Reflection produces polylogarithms of order one}
For $0<a<1$, Gamma reflection gives
\[
 \Rea\log\Gamma(a+iy)+\Rea\log\Gamma(1-a+iy)
  =\log(2\pi)-\tfrac12\log\bigl(2(\cosh(2\pi y)-\cos(2\pi a))\bigr).
\]
Consequently, with $z_t=\exp(2\pi ia-2\pi bt)$,
\begin{equation}\label{hsd:eq:reflected-atan}
 S_a(b)+S_{1-a}(b)=2\pi b\log2+
 2\Rea\int_0^1\frac{\Li_1(z_t)-\Li_1(z_1)}{1-t^2}\dd t.
\end{equation}
To verify the constant, write the logarithm as
$2\pi bt-2\Rea\Li_1(z_t)$ and use
$\int_0^1(1-t)/(1-t^2)\dd t=\log2$.
The integral is ordinary and has a removable endpoint at one.
At $a=1/2$, this yields
\begin{equation}\label{hsd:eq:half-atan}
 S_{1/2}(b)=\pi b\log2+
 \int_0^1\frac{\Li_1(-e^{-2\pi bt})-\Li_1(-e^{-2\pi b})}{1-t^2}\dd t.
\end{equation}
Taking $a\downarrow0$ in \eqref{hsd:eq:reflected-atan} gives
\begin{equation}\label{hsd:eq:integer-atan}
 S_1(b)=-\frac{\pi^2}{8}+\pi b\log2+
 \int_0^1\frac{\Li_1(e^{-2\pi bt})-\Li_1(e^{-2\pi b})}{1-t^2}\dd t.
\end{equation}
Here the logarithmic singularity at $t=0$ is integrable. On the series
side $S_a(b)\to\pi^2/4+S_1(b)$, which accounts for the otherwise easily
missed term $-\pi^2/8$.

\section{Second jets at every nonpositive integer: finite polylogarithmic integrands}
\label{hsd:sec:negative}
Let $0<a<1$, $b>0$ and
\begin{equation}\label{hsd:eq:E}
 E(s;a,b)=\sum_{n\in\mathbb Z}\bigl((n+a)^2+b^2\bigr)^{-s}
         =Z_1(s;a,b)+Z_1(s;1-a,b).
\end{equation}
Put $L_m(a,b)=E'(-m;a,b)$ and $M_m(a,b)=E''(-m;a,b)$.
The first-jet formula of the predecessor \cite{hsd:predecessor} is
\begin{align}
 L_m(a,b)={}&-\frac{2^{2m+1}(m!)^2\pi}{(2m+1)!}b^{2m+1}\notag\\
 &+2(-1)^m m!\sum_{j=0}^m
 \frac{(m+j)!b^{m-j}}{j!(m-j)!4^j\pi^{m+j}}
             \Rea\Li_{m+j+1}(e^{-2\pi b+2\pi ia}).
 \label{hsd:eq:Lfinite}
\end{align}
In particular,
\begin{align}
 L_0(a,b)&=-\log\bigl(2(\cosh(2\pi b)-\cos(2\pi a))\bigr),
 \label{hsd:eq:L0}\\
 L_1(a,b)&=-\frac{4\pi b^3}{3}-\frac{2b}{\pi}\Rea\Li_2(z)
                              -\frac1{\pi^2}\Rea\Li_3(z).
 \label{hsd:eq:L1}
\end{align}
These inherited formulas can be obtained from Poisson summation and the
finite half-integer $K$-Bessel polynomial. They also satisfy
$\partial_{b^2}L_m=mL_{m-1}$, with the zero-mass values fixed by the
Hurwitz functional equation. The new second-jet theorem below uses their
\emph{finite explicit right sides}, not an unevaluated first-jet symbol.

\begin{theorem}[All negative-point second jets]\label{hsd:thm:negative-second}
For every $m\ge0$, $0<a<1$ and $b>0$, define
\begin{equation}\label{hsd:eq:Rintegrand}
 R_m(t;a,b)=\sum_{\ell=0}^m\binom m\ell
       \bigl(b^2(1-t^2)\bigr)^{m-\ell}L_\ell(a,bt).
\end{equation}
Then
\begin{align}
 M_m(a,b)={}&4\Rea\sum_{\alpha\in\{a,1-a\}}\sum_{j=0}^m
 \binom mj(-2ib)^{m-j}\zeta''(-m-j,\alpha+ib)\notag\\
 &-4\int_0^1\frac{L_m(a,b)-R_m(t;a,b)}{1-t^2}\dd t.
 \label{hsd:eq:Mfinite}
\end{align}
The integrand is explicit in finitely many ordinary polylogarithms via
\eqref{hsd:eq:Lfinite}; the endpoint at $t=1$ is removable and the
endpoint at $t=0$ is regular. No Bessel-order derivative or unspecified
finite part occurs.
\end{theorem}
\begin{proof}
Apply Theorem~\ref{hsd:thm:gamma-second} at each of the two shifts with
the fixed even polynomial $W_b(x)=(x^2+b^2)^m$. Holding this polynomial
fixed while differentiating the spectral variable gives
$Z_{W_b}''(0;\alpha,b)=Z_1''(-m;\alpha,b)$ and $C_{W_b}=0$.
If $y=x+ib$, then
\[
 (x^2+b^2)^m=y^m(y-2ib)^m
 =\sum_{j=0}^m\binom mj(-2ib)^{m-j}y^{m+j}.
\]
Therefore the linear endpoint contribution is exactly the finite sum in
\eqref{hsd:eq:Mfinite}.

For an arbitrary fixed even polynomial $W$, the first defect vanishes,
so
\[
 T_W(a,y):=Z_W'(0;a,y)+Z_W'(0;1-a,y)
 =2\Rea\{\Phi_{a,W}(iy)+\Phi_{1-a,W}(iy)\}.
\]
Equation~\eqref{hsd:eq:Kreal} consequently gives
\begin{equation}\label{hsd:eq:bilatK}
 \KK_{a,W}(ib,-ib)+\KK_{1-a,W}(ib,-ib)
 =4\int_0^1\frac{T_W(a,b)-T_W(a,bt)}{1-t^2}\dd t.
\end{equation}
For $W=W_b$, expand instead around the intermediate mass $bt$:
\[
 (x^2+b^2)^m=\sum_{\ell=0}^m\binom m\ell
      \bigl(b^2(1-t^2)\bigr)^{m-\ell}(x^2+b^2t^2)^\ell.
\]
This proves $T_{W_b}(a,bt)=R_m(t;a,b)$ and
$T_{W_b}(a,b)=L_m(a,b)$, hence \eqref{hsd:eq:Mfinite}.

At $t=1$, $R_m(1)=L_m$ and all functions are differentiable, so the
quotient has a finite limit. At zero the bilateral lattice has a positive
gap because $0<a<1$; its continued jets are analytic in $b^2$ near zero.
Equivalently, the polylogarithms at $e^{2\pi ia}\ne1$ in
\eqref{hsd:eq:Lfinite} have regular limits. This proves the endpoint claims.
\end{proof}

For $m=0$ this is the transparent formula
\begin{align}
 E''(0;a,b)={}&4\Rea\{\zeta''(0,a+ib)+\zeta''(0,1-a+ib)\}
                         +4S_a(b)+4S_{1-a}(b)\notag\\
 ={}&4\Rea\{\zeta''(0,a+ib)+\zeta''(0,1-a+ib)\}
 +8\pi b\log2\notag\\
 &+8\Rea\int_0^1\frac{\Li_1(e^{2\pi ia-2\pi bt})
                      -\Li_1(e^{2\pi ia-2\pi b})}{1-t^2}\dd t.
 \label{hsd:eq:M0explicit}
\end{align}
This is a representation theorem, not a nonreducibility theorem. In
particular, retaining an ordinary polylogarithmic integral is not evidence
that no further finite reduction exists.

\subsection{A normalized all-jet mass-antiderivative ladder}
The second jets themselves satisfy an inhomogeneous relation:
\begin{equation}\label{hsd:eq:M-inhomogeneous}
 \partial_{\mu}M_m(a,\sqrt\mu)
       =mM_{m-1}(a,\sqrt\mu)-2L_{m-1}(a,\sqrt\mu),\qquad m\ge1.
\end{equation}
Harmonic normalization removes the inhomogeneous term. In fact the same
normalization works at every spectral order.

Let $h_{m,q}$ be the complete homogeneous polynomial in
$1,1/2,\ldots,1/m$, defined by
\begin{equation}\label{hsd:eq:completeH}
 \sum_{q\ge0}h_{m,q}t^q=\prod_{j=1}^m(1-t/j)^{-1}
 =\exp\left(\sum_{p\ge1}\frac{H_m^{(p)}}p t^p\right).
\end{equation}
For $m=0$ the product is one. Define, for $r\ge1$,
\begin{equation}\label{hsd:eq:normalizedjet}
 \mathcal A_{m,r}(a,\mu)=
 \left.\partial_t^r\left\{\frac{m!}{(1-t)_m}
                        E(t-m;a,\sqrt\mu)\right\}\right|_{t=0}.
\end{equation}

\begin{theorem}[Every normalized jet is an Appell-type mass ladder]
\label{hsd:thm:Appell}
For $0<a<1$, $\mu\ge0$, $m\ge1$ and $r\ge1$,
\begin{align}
 \mathcal A_{m,r}(a,\mu)&=r!\sum_{j=1}^r
                 \frac{h_{m,r-j}}{j!}E^{(j)}(-m;a,\sqrt\mu),
                 \label{hsd:eq:A-explicit}\\
 \partial_\mu\mathcal A_{m,r}(a,\mu)&=m\mathcal A_{m-1,r}(a,\mu).
                 \label{hsd:eq:A-ladder}
\end{align}
Consequently every anchored mass primitive is explicit:
\begin{align}
 \mathcal A_{m,r}(a,\mu)={}&\sum_{j=0}^{m-1}\binom mj\mu^j
                                  \mathcal A_{m-j,r}(a,0)\notag\\
 &+m\int_0^\mu(\mu-u)^{m-1}E^{(r)}(0;a,\sqrt u)\dd u.
 \label{hsd:eq:A-primitive}
\end{align}
The boundary constants are finite Hurwitz jets because
\begin{equation}\label{hsd:eq:masszerojets}
 E^{(j)}(-m;a,0)=2^j\{\zeta^{(j)}(-2m,a)+\zeta^{(j)}(-2m,1-a)\}.
\end{equation}
In particular,
$\mathcal A_{m,1}=L_m$ and
$\mathcal A_{m,2}=M_m+2H_mL_m$.
\end{theorem}
\begin{proof}
The parameter identity $\partial_\mu E(s)=-sE(s+1)$ holds first in
absolute convergence and then by continuation. Therefore, as germs in $t$,
\[
 \partial_\mu\left\{\frac{m!}{(1-t)_m}E(t-m)\right\}
 =m\frac{(m-1)!}{(1-t)_{m-1}}E(t-(m-1)).
\]
Differentiate in $t$ to obtain \eqref{hsd:eq:A-ladder}.
The identities $E(-m)=0$ follow from the binomial continuation and
$\zeta(-2h,a)+\zeta(-2h,1-a)=0$, the latter being Bernoulli reflection.
Expanding the prefactor by \eqref{hsd:eq:completeH} now gives
\eqref{hsd:eq:A-explicit}. The positive lattice gap permits continuation
and parameter differentiation through $\mu=0$; at zero the identity
$E(s;a,0)=\zeta(2s,a)+\zeta(2s,1-a)$ proves
\eqref{hsd:eq:masszerojets}. Repeated integration of the ladder, with
these boundary data, gives \eqref{hsd:eq:A-primitive}.
\end{proof}

Combining \eqref{hsd:eq:A-primitive} at $r=2$ with
Theorem~\ref{hsd:thm:negative-second} evaluates every polynomially weighted
mass integral of $E''(0)$ by explicit endpoints and a single finite
polylogarithmic-integrand coordinate. No integration constants are
chosen after the fact.

\section{Several quadratic factors and the complete confluent expansion}
\label{hsd:sec:masses}
Throughout this section $W(x)=P(x^2)$ is even; its coefficients may be
complex. For $b_1,\ldots,b_r>0$ define
\begin{equation}\label{hsd:eq:manyZ}
 Z_{W;\mathbf b}(s;a)=\sum_{n\ge0}W(x_n)
                  \prod_{j=1}^r(x_n^2+b_j^2)^{-s}.
\end{equation}
The same notation with a superscript $\mathrm{bil}$ denotes the sum over
the full lattice, equivalently the sum of the one-sided functions at
$a$ and $1-a$, when $0<a<1$.

For $\mu>0$ put
\begin{equation}\label{hsd:eq:Tmu}
 \mathcal T_{a,W}(\mu)=Z_W'(0;a,\sqrt\mu)
        =\Phi_{a,W}(i\sqrt\mu)+\Phi_{a,W}(-i\sqrt\mu).
\end{equation}
The last equality is the inherited first-defect parity formula.
On the bilateral lattice it is the following finite polylogarithmic
expression:
\begin{equation}\label{hsd:eq:Tbil}
 \mathcal T_{a,W}^{\mathrm{bil}}(\mu)
       =\sum_{\ell=0}^{\deg P}\frac{P^{(\ell)}(-\mu)}{\ell!}
                                  L_\ell(a,\sqrt\mu).
\end{equation}
Thus the next definition uses already evaluated first jets only.

\begin{definition}[Mass covariance]\label{hsd:def:V}
For $\mu=b^2$, $\nu=c^2$, define
\begin{equation}\label{hsd:eq:Vmass}
 V_{a,W}(b,c)=\int_0^1
 \frac{\mathcal T_{a,W}(\mu)+\mathcal T_{a,W}(\nu)
 -\mathcal T_{a,W}((1-t)\mu+t\nu)
 -\mathcal T_{a,W}(t\mu+(1-t)\nu)}{t(1-t)}\dd t.
\end{equation}
Use $\mathcal T^{\mathrm{bil}}$ for $V^{\mathrm{bil}}$.
The name denotes this explicitly defined analytic function; it does not
imply a probability interpretation for arbitrary polynomial weights.
\end{definition}

The integrand has removable endpoints and is identically zero when
$b=c$. The formula is analytic in $\mu,\nu$ on a neighborhood of every
positive real rectangle, including the diagonal.

\begin{theorem}[All second jets are pairwise mass covariances]
\label{hsd:thm:masses}
For every even polynomial $W$, every $a>0$, and all positive masses,
\begin{equation}\label{hsd:eq:many-second}
 \boxed{\displaystyle
 Z_{W;\mathbf b}''(0;a)=r\sum_{j=1}^r Z_W''(0;a,b_j)
                    -\sum_{1\le i<j\le r}V_{a,W}(b_i,b_j).}
\end{equation}
In particular,
\begin{equation}\label{hsd:eq:two-second}
 Z_{W;b,c}''(0;a)=2Z_W''(0;a,b)+2Z_W''(0;a,c)-V_{a,W}(b,c).
\end{equation}
Every term is explicit by Theorem~\ref{hsd:thm:gamma-second} and
\eqref{hsd:eq:Vmass}. The bilateral statements follow by reflection and
have finite polylogarithmic covariance integrands via
\eqref{hsd:eq:Tbil}. Repeated masses are allowed; in particular
$Z_{W;b,b}''(0;a)=4Z_W''(0;a,b)$.
\end{theorem}
\begin{proof}
We first prove the two-factor statement directly from a beta integral,
without a formal manipulation of divergent logarithmic sums. Euler's
identity gives
\[
 Z_{W;b,c}(s;a)=\frac{\Gamma(2s)}{\Gamma(s)^2}
 \int_0^1t^{s-1}(1-t)^{s-1}
 Z_W\bigl(2s;a,\sqrt{(1-t)\mu+t\nu}\bigr)\dd t
\]
in its initial convergence half-plane. Subtract its endpoint average
as in Theorem~\ref{hsd:thm:beta}. For even $W$,
\begin{equation}\label{hsd:eq:quadraticzero}
 Z_W(0;a,b)=\sum_{d\ \mathrm{even}}w_d\zeta(-d,a)
\end{equation}
is independent of $b$: in the inverse-power expansion every nonconstant
Pochhammer coefficient vanishes at zero, and $2k-d$ cannot equal one.
Thus the subtracted zero jet is identically zero. The second beta row
contains only the first-jet chord, exactly \eqref{hsd:eq:Vmass}.
This proves \eqref{hsd:eq:two-second}, including continuation and endpoint
regularity.

For completeness, a common coefficient argument justifies the extension
to all $r$. Choose a tail $x\ge X>\max b_j$ and expand each product as
\[
 x^{-2rs}\sum_{k\ge0}c_k(s)x^{-2k}.
\]
After multiplication by an even polynomial $W$, continuation replaces
each power by $\zeta(2rs+2k-d,X)$. None of these Hurwitz terms has a pole
at $s=0$, since $2k-d$ is even. Therefore taking two spectral derivatives
is equivalent to applying the same linear coefficient rule
\[
 x^e(\log x)^q\longmapsto(-1)^q\zeta^{(q)}(-e,X),
 \qquad e\ \hbox{even},
\]
to the pointwise second-derivative expansion. This rule is independent
of the initial scale $2r$; that scale is already accounted for by the
powers of $\log x$. The expansions converge normally after the finite
head, as in Theorem~\ref{hsd:thm:stieltjes}.

Apply this common linear rule to the elementary identity
\[
 \left(\sum_{j=1}^r A_j\right)^2
   =r\sum_{j=1}^r A_j^2-\sum_{i<j}(A_i-A_j)^2,
 \qquad A_j=\log(x^2+b_j^2).
\]
The two-factor case has already identified the image of
$(A_i-A_j)^2$ as $V_{a,W}(b_i,b_j)$. The finite head satisfies the same
pointwise identity. This proves \eqref{hsd:eq:many-second} without
assuming linearity across a hidden pole. Reflection proves the bilateral
case. The coincident statement follows either from $V(b,b)=0$ or from
$Z_{W;b,b}(s)=Z_W(2s;b)$.
\end{proof}

\subsection{An independent Gamma polarization of the covariance}
There is a second useful expression for the same function:
\begin{align}
 V_{a,W}(b,c)={}&\sum_{\epsilon,\eta\in\{1,-1\}}
                  \KK_{a,W}(\epsilon ib,\eta ic)\notag\\
               &-\KK_{a,W}(ib,-ib)-\KK_{a,W}(ic,-ic).
 \label{hsd:eq:V-Gamma}
\end{align}
To prove it, first use a tail with both masses small relative to $x$.
The algebraic combination of logarithm squares on the right is
$\log^2((x^2+b^2)/(x^2+c^2))$. In the expansion
\eqref{hsd:eq:PhiTaylor}, the extra harmonic term at degree $k$ is
proportional to $w_{k-1}$. The combined logarithm-square expression has
only even inverse powers, whereas even $W$ can contribute to
$w_{k-1}$ only for odd $k$. Hence all such harmonic corrections cancel.
The common coefficient argument in the preceding proof identifies the
remaining expression with $V$. Analytic continuation in positive masses
proves the global identity. Formula~\eqref{hsd:eq:V-Gamma} is not needed
to define $V$ and is valuable as an independent evaluation route.

When $\deg W\le2$, the covariance is the ordinary convergent series
\begin{equation}\label{hsd:eq:Vordinary}
 V_{a,W}(b,c)=\sum_{n\ge0}W(x_n)
                   \log^2\left(\frac{x_n^2+b^2}{x_n^2+c^2}\right),
 \qquad W\ \hbox{even},\quad\deg W\le2.
\end{equation}
Indeed, the logarithm square is $O(x^{-4})$. For higher polynomial
degree the unmodified series generally diverges, and
\eqref{hsd:eq:Vmass}, not that divergent series, is the assertion.

For the unweighted bilateral lattice, define
$D(a,y)=2(\cosh(2\pi y)-\cos(2\pi a))$. The following ordinary identity
is an explicit elementary-integrand specialization:
\begin{align}
 &\sum_{n\in\mathbb Z}
  \log^2\left(\frac{(n+a)^2+b^2}{(n+a)^2+c^2}\right)\notag\\
 &\quad=\int_0^1\frac1{t(1-t)}
 \log\left(\frac{D(a,\sqrt{(1-t)b^2+tc^2})
                  D(a,\sqrt{tb^2+(1-t)c^2})}{D(a,b)D(a,c)}\right)\dd t.
 \label{hsd:eq:V-elementary}
\end{align}
All quantities are positive inside the logarithms, so no branch convention
is implicit.

\subsection{Every confluent coefficient is a finite polygamma expression}
\begin{theorem}[Complete coincident-mass expansion]\label{hsd:thm:collision}
Let $a,b>0$ and $W$ be even. For $h$ sufficiently small, more precisely
$|h|<a^2+b^2$ in the one-sided analytic squared-mass continuation,
\begin{equation}\label{hsd:eq:Vcollision}
 \boxed{\displaystyle
 V_{a,W}(b,\sqrt{b^2+h})=
 \sum_{k\ge2}\frac{2(-1)^kH_{k-1}}k\,h^k Z_W(k;a,b).}
\end{equation}
The left side denotes analytic continuation in $h$ of its positive-mass
germ. The square-root notation creates no singularity at $h=-b^2$,
because the spectrum is even in the mass. The series is locally normally
convergent. On the bilateral lattice, replace the radius by $\min(a,1-a)^2+b^2$ and use the bilateral zeta.
All displayed positive-integer values are holomorphic continued values,
and each has a finite polygamma--Bernoulli evaluation. The first terms are
\begin{equation}\label{hsd:eq:Vfirstterms}
 h^2Z_W(2;a,b)-h^3Z_W(3;a,b)+\frac{11}{12}h^4Z_W(4;a,b)+O(h^5).
\end{equation}
\end{theorem}
\begin{proof}
The even-weight continuation has no positive-integer spectral poles:
its possible poles are at half-integers, as is immediate from the
inverse-power expansion. Repeated mass differentiation of the first
jet therefore gives
\begin{equation}\label{hsd:eq:Tderivs}
 \mathcal T_{a,W}^{(k)}(\mu)=(-1)^k(k-1)!Z_W(k;a,\sqrt\mu),
 \qquad k\ge1.
\end{equation}
Expand the ordinary analytic function $\mathcal T$ in
\eqref{hsd:eq:Vmass} about $\mu=b^2$. The coefficient of $h^k$ is
\[
 \frac{\mathcal T^{(k)}(\mu)}{k!}
 \int_0^1\frac{1-t^k-(1-t)^k}{t(1-t)}\dd t
       =\frac{2H_{k-1}}{k!}\mathcal T^{(k)}(\mu).
\]
The elementary integral follows by finite geometric summation. This
proves the coefficients. The nearest zero of $x_n^2+\mu+h$ is at
$h=-a^2-b^2$; the finite-head continuation is analytic in the open disk
bounded by that point. The Taylor expansion and the endpoint subtraction
are locally uniform on smaller disks, proving normal convergence.
The bilateral radius follows from the smallest lattice gap.
\end{proof}

Here is the promised finite evaluation of the coefficients. For $k\ge1$,
\begin{align}
 Z_1(1;a,b)&=\frac{\psi(a+ib)-\psi(a-ib)}{2ib},\label{hsd:eq:resolvent}\\
 Z_1(k;a,b)&=\frac{(-1)^{k-1}}{(k-1)!}
   \left(\frac{\partial}{\partial\mu}\right)^{k-1}
   \frac{\psi(a+i\sqrt\mu)-\psi(a-i\sqrt\mu)}{2i\sqrt\mu}\bigg|_{\mu=b^2}.
   \label{hsd:eq:positivepoly}
\end{align}
Only finitely many ordinary polygamma derivatives occur. For $m\ge0$,
\begin{equation}\label{hsd:eq:nonpositivepoly}
 Z_1(-m;a,b)=\sum_{j=0}^m\binom mj b^{2(m-j)}\zeta(-2j,a),
\end{equation}
whose Hurwitz values are Bernoulli polynomials. Expanding the weight
around $-b^2$ gives, for every positive integer $k$,
\begin{equation}\label{hsd:eq:weightedresolvent}
 Z_W(k;a,b)=\sum_{\ell=0}^{\deg P}\frac{P^{(\ell)}(-b^2)}{\ell!}
                                        Z_1(k-\ell;a,b).
\end{equation}
Equations~\eqref{hsd:eq:positivepoly} and
\eqref{hsd:eq:nonpositivepoly} cover all terms. No confluent denominator
needs to be evaluated by subtracting nearby masses.

A related mixed derivative identity is
\begin{equation}\label{hsd:eq:mixedmass}
 \partial_\mu\partial_\nu Z_{W;\sqrt\mu,\sqrt\nu}''(0;a)
   =2\frac{Z_W(1;a,\sqrt\mu)-Z_W(1;a,\sqrt\nu)}{\nu-\mu},
\end{equation}
with diagonal value $2Z_W(2;a,\sqrt\mu)$. To prove it, differentiate the
initial spectral series twice in the masses, obtaining a factor $s^2$,
and then take the second spectral derivative. Rational partial fractions
give the right side. The same even-power pole exclusion used above
justifies this continuation. Thus both the covariance and its mass
Hessian have completely specified, regular diagonal limits.

\section{Anchored shift primitives of every spectral order}
\label{hsd:sec:primitives}
The preceding representations also admit a uniform antiderivative
calculus in the Hurwitz shift $a$. This is separate from the mass ladder:
the two parameters should not be interchanged informally.

For integers $R\ge1$, $h,j\ge0$, define
\begin{equation}\label{hsd:eq:Pgeneral}
 \PP_{R,h}^{[j]}(a)=(-1)^R
 \left.\partial_s^j\left(\frac{\zeta(s-R,a)}{(s-R)_R}\right)\right|_{s=-h}.
\end{equation}
The denominator has no zero at $s=-h$; all its factors are strictly
negative integers there. Thus the formula is a finite combination of
negative-integer Hurwitz derivatives and generalized harmonic numbers,
with no resonance prescription needed. Its basic property is
\begin{equation}\label{hsd:eq:Pderivative}
 \partial_a^R\PP_{R,h}^{[j]}(a)=\zeta^{(j)}(-h,a).
\end{equation}
For the Gamma case $j=1$, this reads explicitly
\begin{equation}\label{hsd:eq:Pexplicit}
 \PP_{R,h}^{[1]}(a)=\frac{h!}{(h+R)!}
 \left\{\zeta'(-h-R,a)+(H_{h+R}-H_h)\zeta(-h-R,a)\right\}.
\end{equation}
For example, $\partial_a\{\zeta'(-1,a)+\zeta(-1,a)\}=\zeta'(0,a)$.
The elementary $\tfrac12\log(2\pi)$ can be added when a primitive of
$\log\Gamma(a)$ rather than $\zeta'(0,a)$ is desired.

To fix all constants, let $a_0>0$ and write
\begin{equation}\label{hsd:eq:anchorop}
 \mathfrak A_{a_0,R}Q(a)=Q(a)-\sum_{q=0}^{R-1}
                    \frac{(a-a_0)^q}{q!}Q^{(q)}(a_0).
\end{equation}
This is the unique $R$th primitive of $Q^{(R)}$ whose derivatives of
orders $0,\ldots,R-1$ vanish at $a_0$.

\begin{theorem}[All anchored shift primitives of the defect]
\label{hsd:thm:shiftprimitive}
For fixed admissible endpoints $u,v$ and polynomial $W$, every
$\Delta_r(a,W;u,v)$ has an explicitly anchored $R$th primitive on any
positive $a$-interval where the endpoint segments remain admissible.
In Theorem~\ref{hsd:thm:beta-jets}, replace each finite endpoint function
\begin{equation}\label{hsd:eq:Fjfinite}
 F_{a,W}^{(j)}(0;z)=\sum_{d=0}^D w_d\sum_{h=0}^d
                 \binom dh(-z)^{d-h}\zeta^{(j)}(-h,a+z)
\end{equation}
by
\begin{equation}\label{hsd:eq:anchoredF}
 \mathfrak A_{a_0,R}\left[
 \sum_{d=0}^D w_d\sum_{h=0}^d\binom dh(-z)^{d-h}
                            \PP_{R,h}^{[j]}(a+z)\right].
\end{equation}
Keep the beta coefficients and logarithmic weights unchanged. The
result is the required anchored primitive, still a finite sum of
ordinary convergent integrals.
\end{theorem}
\begin{proof}
The classical parameter identity
$\partial_a^R\zeta(s-R,a)=(-1)^R(s-R)_R\zeta(s,a)$
proves \eqref{hsd:eq:Pderivative} by differentiation in $s$.
Differentiating the reciprocal finite product at $s=-h$ gives
\eqref{hsd:eq:Pexplicit}. The operator
\eqref{hsd:eq:anchorop} imposes the stated initial conditions without
changing the $R$th derivative. The endpoint subtraction commutes with
this operation because it is linear and the endpoints are fixed.
On compact admissible $a$-intervals the endpoint estimates in
Theorem~\ref{hsd:thm:beta} are uniform, including the fixed logarithmic
weights. Therefore the $R$ shift derivatives may pass under every
integral. The resulting $R$th derivative is precisely
\eqref{hsd:eq:all-beta}, and all lower derivatives at the anchor vanish.
\end{proof}

In the conjugate second-order case there is a shorter equivalent formula.
Construct $\Phi_{a,W}^{\langle R\rangle}(z)$ by replacing the first jets
in \eqref{hsd:eq:phi} with the anchored version of
\eqref{hsd:eq:Pexplicit}. Then
\begin{equation}\label{hsd:eq:anchoredDelta2}
 \int_{a_0}^{a}\frac{(a-x)^{R-1}}{(R-1)!}\Delta_2(W;x,b)\dd x
 =\int_0^1\frac{\CC_t\Phi_{a,W}^{\langle R\rangle}(ib,-ib)}{t(1-t)}\dd t
                 -C_W(b)\frac{(a-a_0)^R}{R!}.
\end{equation}
The formula also holds for $a<a_0$ with oriented integration.
It explicitly retains the primitive of the local polynomial, which must
not be dropped merely because that polynomial is independent of $a$.

\subsection{Parameter differentiation and generalized Stieltjes functions}
The coefficient formula \eqref{hsd:eq:Stieltjes-master} is locally normal
in $a,b$ after a sufficiently large fixed head is selected on a compact
parameter set. It can therefore be differentiated in either parameter
any prescribed finite number of times. At resonant coordinates use
\eqref{hsd:eq:Stieltjes} as a whole germ, rather than differentiating a
singular isolated zeta term. In particular,
\[
 \partial_a\zeta(s,a)=-s\zeta(s+1,a)
\]
converts the resulting derivatives into polygamma and Stieltjes/Hurwitz
jets with explicit finite product coefficients. At an integer shift,
$\zeta(k,a)$ is the corresponding generalized harmonic tail for $k\ge2$.
Thus the finite beta representation and the harmonic/Stieltjes series
remain equal after all these legitimate differentiations and after the
anchored integrations just proved.

\section{Verification, correction guards, and integration status}
\label{hsd:sec:audit}
\subsection{What was actually checked}
The mathematical theorems above are established by their analytic proofs.
The accompanying computations check finite algebra and implementations;
none is substituted for an all-parameter proof.

The exact suite passed \textbf{432 assertions} over rational polynomial
algebra and formal Stieltjes coefficients. Its groups are:
\begin{center}
\begin{tabular}{lr}
\toprule
Finite algebra checked & Assertions\\
\midrule
Double-zero and duplication polynomials &60\\
Harmonic coefficients and Gamma corrections &60\\
Local beta moments through the required logarithmic order &32\\
Confluent log-square coefficients and beta moments &62\\
Shift-primitive normalization &128\\
All-jet mass-germ normalization &48\\
Explicit Stieltjes pole-jet extraction &42\\
\midrule
Total &432\\
\bottomrule
\end{tabular}
\end{center}
The tests include $D_k$ through $k=20$, the two local beta moments through
$k=16$, the confluent coefficient identity through $k=32$, anchored
primitive coefficients for $1\le R\le8$ and $0\le h\le7$, and the mass
germ identity through $m=24$. The complete finite assertions are recorded
in \src{results/exact_results.json}.

The numerical suite completed \textbf{87 comparisons at 35 decimal-digit
working precision}. The acceptance threshold was $10^{-23}$ for
$|\mathrm{lhs}-\mathrm{rhs}|/\max(1,|\mathrm{lhs}|,|\mathrm{rhs}|)$.
The largest recorded scaled residual was approximately
$9.18\times10^{-35}$. This is a residual, not an error bound.
The checks cover all of the following mechanisms: independent direct
Pochhammer and double-zero defect jets through spectral order four;
odd and even Gamma corrections; cutoff changes and masses outside the
unshifted Taylor disk; ordinary logarithm and arctangent sums; reflected
polylogarithmic integrals including a logarithmic endpoint; second jets
at four nonpositive spectral points and two shift/mass pairs; two distinct
Gamma/mass covariance representations; independently expanded two-factor
spectra; repeated masses; finite polygamma confluence; shift primitives;
and logarithmically weighted beta jets.

Some comparisons share the same underlying numerical Hurwitz-zeta
implementation. They must not be advertised as entirely independent
special-function libraries. The ordinary sums, Gamma products,
polylogarithmic integrands and spectral expansions nevertheless provide
different algebraic and analytic evaluation routes. No interval
enclosures, certified infinite tails, or proof-assistant verification
are included. Execution was completed in staged invocations; the final
JSON contains every accepted comparison, not an extrapolated count.

\subsection{Corrections versus safeguards}
No false theorem was identified in the inspected predecessor's
first-derivative statements. In particular it explicitly cautions that
higher derivatives do not automatically factorize. The following are
therefore proposed \emph{integration safeguards}, not accusations that
the predecessor made those errors.

\paragraph{Do not extend first-order parity to the full second jet.}
Even weights remove $C_W$ and the $I_{0,p}$ terms, but not $\KK$.
The exact counterexample is $\Delta_2(1;a,b)=-4S_a(b)<0$ for $b>0$.
Theorem~\ref{hsd:thm:nonvanishing} gives the full universal criterion.

\paragraph{Retain one extra coefficient before a pole multiplication.}
The resonant $r$th derivative needs $d_{k,r+1}$, not only $d_{k,r}$.
Equation~\eqref{hsd:eq:Tres} is the exact replacement rule. This is also
why generalized Stieltjes coefficients and ordinary zeta derivatives
must not be mixed without tracking the spectral scale two.

\paragraph{Do not silently remove the Gamma correction for odd weights.}
For $W=x$ the second defect is
$\Delta_2(x;a,b)=\KK_{a,x}(ib,-ib)+4b^2$.
Omitting the $4b^2$ term is an explicitly false identity.

\paragraph{Normalize the second mass ladder.}
The pure ladder belongs to $M_m+2H_mL_m$, not to $M_m$ alone.
At all orders the correct object is the germ
$m!E(t-m)/(1-t)_m$. The primitive constants are those in
\eqref{hsd:eq:masszerojets}, not arbitrary polynomials.

\paragraph{Distinguish continuation from an ordinary sum.}
The covariance series in \eqref{hsd:eq:Vordinary} is ordinary for even
weights of degree at most two. At higher degree use the ordinary
\emph{integral} definition \eqref{hsd:eq:Vmass}; the original logarithmic
series is generally divergent. Similarly, a small-mass series is not
globally valid without its finite head.

\paragraph{Preserve the branch and endpoint subtractions.}
Use analytic $\log\Gamma$, retain the paired numerator under the
integral, and evaluate the diagonal by its confluent formula. Integrating
its four separated divergent endpoint terms is not equivalent to the
subtracted beta integral.

\subsection{Suggested repository integration}
A suitable new report destination is
\begin{quote}\small
\src{Analysis/Polylogarithms/docs/reports/stieltjes-correlation-calculus/continuations/higher-spectral-defects/}.
\end{quote}
The package includes a compact theorem fragment with the prefix
\src{hsd:}, source hashes, a machine-readable claim ledger, and an
editorial status note. The predecessor's Section~7.6 Questions~2 and~7
can be cross-referenced to the explicit representation theorems here.
Question~1 can record the new finite-integrand representation while
retaining its minimum-function-algebra component as open. The classical
first-derivative results should keep their original attribution. No
Gaussian special-value status or numerical independence status should
be changed. No remote commit, overwrite, or repository modification
was performed in producing this package.

\section{Further research questions and concrete next targets}
\label{hsd:sec:future}
\paragraph{1. Finite closure of the negative-point second integrals.}
Determine special parameter families for which the integral in
\eqref{hsd:eq:Mfinite} reduces to a finite collection of polylogarithms,
polylogarithm order derivatives, and explicit logarithms. A useful first
target is the half-shift $a=1/2$ at $m=0,1$. A failure in one prescribed
symbolic calculus is not a proof of analytic nonreducibility.

\paragraph{2. A finite polylogarithmic dictionary for every higher jet.}
Theorem~\ref{hsd:thm:beta-jets} lowers the inner Hurwitz order by one.
Can reflection and Fourier continuation turn those lower jets into a
uniform finite set of logarithmically weighted polylogarithm-order
integrals at every $r$? The goal should include the exact endpoint
constants and not merely a formal interchange of a Fourier series with
a singular beta weight.

\paragraph{3. Higher-order multi-mass connected terms.}
The pairwise formula \eqref{hsd:eq:many-second} relies on a quadratic
variance identity. At third and higher order, identify the appropriate
three-point and higher endpoint subtractions in a Dirichlet-simplex
parameterization. Determine which terms genuinely require more than
pairwise covariances and how their coincident limits encode products of
harmonic coefficients and higher resolvents.

\paragraph{4. Multi-mass formulas for arbitrary odd polynomial weights.}
For odd weights the zero quadratic jet depends polynomially on the mass,
so the local term in the mass beta formula need not vanish. Compute its
all-order residue polynomial and the corresponding corrected multi-mass
covariance formula. The all-order linear-factor beta master and the
explicit $J_k$ moments provide exact interpolation constraints.

\paragraph{5. Arithmetic and colored weights.}
Replace the polynomial weight by a polynomial times a periodic sequence.
Separate its mean, which controls the arithmetic Hurwitz pole, from the
nontrivial character components. Seek a small-denominator Gamma/Dirichlet
$L$ dictionary for the resulting beta polarizations and their Stieltjes
jets, while keeping analytic evaluation separate from claims of period
independence.

\paragraph{6. Joint lattice-gap and mass degeneration.}
The formulas are regular for $0<a<1$ and positive mass, but the corner
$(a,b)=(0,0)$ has a collapsing lattice point. Construct a two-parameter
subtracted germ that separates that point explicitly. Determine which
normalized mixed jets have path-independent limits and which retain
specified logarithmic path terms.

\paragraph{7. Nonintegral polynomial degree.}
For weights $x^\rho$, the finite shifted-Hurwitz expression
\eqref{hsd:eq:Ffinite} is no longer available. Develop an analytic
replacement that interpolates the odd-degree residue corrections as
$\rho$ crosses integers. A correct interpolation must reproduce both
$\Delta_1$ and $C_W$ at all polynomial weights, rather than assuming the
second correction is a continuation of first-order parity.

\paragraph{8. Formalization and validated evaluation.}
The beta subtraction, polynomial pole extraction, harmonic recurrences,
and normalized primitive identities have clean finite algebraic layers.
Formalize those first, then the normal-convergence and differentiation
lemmas. Independently, equip the ordinary beta/Gamma integrals and
finite-head tails with interval error control. Such a validated evaluator
would strengthen the supplied diagnostics without changing the logical
role of the analytic proofs.

\section*{Conclusion}
\addcontentsline{toc}{section}{Conclusion}
The second spectral derivative is not governed by the first-derivative
parity law. Its nonlocal part is nevertheless explicit: a Gamma--Hurwitz
polarization, with a finite harmonic correction for odd weights.
Endpoint-subtracted beta parameterization extends this mechanism to
every spectral order and lowers the inner order by one. The independent
harmonic--Stieltjes series supplies every resonant coefficient and a
globally valid finite-head implementation.

The same construction controls several quadratic masses, gives regular
confluent polygamma expansions, and converts the predecessor's finite
first-jet polylogarithms into explicit second-jet integrands. Harmonic
normalization then produces an all-jet mass-antiderivative ladder, while
negative Hurwitz derivatives fix every shift-primitive constant.
These are exact representation and structural theorems. The remaining
questions concern further finite closure, arithmetic specialization,
higher connected terms, and formal or interval certification, not
unreported conjectural steps in the identities proved here.

\clearpage
\begin{thebibliography}{9}
\raggedright
\bibitem{hsd:repo}
V. Reshetnikov and ProveIt contributors,
\emph{ProveIt: Polylogarithms manuscript}, canonical preface,
differentiation chapter, research-status chapter, and incoming-report
catalogue, inspected 11 October 2026. Repository checkpoint
\pin. \url{https://github.com/VladimirReshetnikov/ProveIt}.
The source audit distinguishes preliminary main-branch reads from the
byte-identified predecessor below.

\bibitem{hsd:predecessor}
\emph{Quadratic-Spectrum Stieltjes Identities: Complete reflected jets,
harmonic counterterms, finite polylogarithmic reductions, and a polynomial
parity criterion}, ProveIt research continuation, 10 October 2026.
Especially Sections~5--7 and Section~7.6 Questions~1, 2, and~7.
Byte-identified single-file source at
\src{Analysis/Polylogarithms/docs/reports/stieltjes-correlation-calculus/continuations/quadratic-spectrum-stieltjes/article_singlefile.tex},
repository checkpoint \pin. SHA-256 recorded in Section~\ref{hsd:sec:scope}.
The predecessor's own source anchor is
\src{446dca9bd65c20e409d1d2561d4ffad408103e2a}.

\bibitem{hsd:intake}
ProveIt, \emph{Sixteenth incoming-archive manifest},
\src{Analysis/Polylogarithms/docs/manuscript/verification/sixteenth-incoming-archives.json},
checkpoint \pin. Consulted for the predecessor's path, byte length, and
SHA-256; not as a substitute for reading its mathematical text.

\bibitem{hsd:beta}
NIST Digital Library of Mathematical Functions, \S5.12,
\emph{Beta function}, especially Euler's integral and Eq.~5.12.4.
\url{https://dlmf.nist.gov/5.12}. Consulted 11 October 2026.
The two-denominator parameterization used here is a direct specialization.

\bibitem{hsd:gamma}
NIST Digital Library of Mathematical Functions, \S5.8,
\emph{Infinite products}; \S5.5, \emph{Functional relations}; and \S5.15,
\emph{Polygamma functions}.
\url{https://dlmf.nist.gov/5.8}; \url{https://dlmf.nist.gov/5.5};
\url{https://dlmf.nist.gov/5.15}.

\bibitem{hsd:hurwitz}
NIST Digital Library of Mathematical Functions, \S25.11,
\emph{Hurwitz zeta function}: continuation, parameter derivatives,
negative-integer values, and the derivative at zero.
\url{https://dlmf.nist.gov/25.11}. Consulted 11 October 2026.

\bibitem{hsd:polylog}
NIST Digital Library of Mathematical Functions, \S25.12,
\emph{Polylogarithms}.
\url{https://dlmf.nist.gov/25.12}. Consulted 11 October 2026.

\bibitem{hsd:dowker}
J. S. Dowker, \emph{Calculation of the multiplicative anomaly},
arXiv:1412.0549, 2014. \url{https://arxiv.org/abs/1412.0549}.
Cited for the classical first-derivative anomaly setting and factorization
scope. The endpoint-subtracted higher-jet identities above are proved
with their own specified normalization.
\end{thebibliography}
\end{document}
